\documentclass{article}
%Load mhchem using some package options
\usepackage[version=4,arrows=pgf-filled,
textfontname=sffamily,
mathfontname=mathsf]{mhchem}
\begin{document}
	\subsection*{Simple molecules}
	
	\begin{itemize}
		\item Water: \verb|\ce{H2O}|, \ce{H2O}
		\item Benzene: \verb|\ce{C6H6}|, \ce{C6H6}
		\item Hydrogen peroxide: \verb|\ce{H2O2}|, \ce{H2O2}
		\item Acetic acid: \verb|\ce{C2H4O2}|, \ce{C2H4O2}
		\item Glucose: \verb|\ce{C6H12O6}|, \ce{C6H12O6}
	\end{itemize}
	
	\subsection*{Chemical equations}
	
	Two basic examples:
	
	\begin{itemize}
		\item \verb|\ce{2H2 + O2 -> 2H2O}| typesets \ce{2H2 + O2 -> 2H2O}
		\item \verb|\ce{CO2 + C -> 2 CO}| typesets \ce{CO2 + C -> 2 CO}
	\end{itemize}
	
	\subsection*{A more complex example}
	
	Writing \verb|\ce{Hg^2+ ->[I-] HgI2 ->[I-] [Hg^{II}I4]^2-}| typesets this:\vskip10pt
	\noindent \ce{Hg^2+ ->[I-] HgI2 ->[I-] [Hg^{II}I4]^2-}.
	
	\subsection*{A math mode example}
	
	Chemical expressions can be typeset using math mode commands such as \verb|\frac|. \vskip10pt
	
	\noindent Writing \verb|\[K=\frac{[\ce{Hg^2+}][\ce{Hg}]}{[\ce{Hg2^2+}]}\]| produces this:
	
	\[K=\frac{[\ce{Hg^2+}][\ce{Hg}]}{[\ce{Hg2^2+}]}\]
\end{document}